\documentclass{article}%
\usepackage[T1]{fontenc}%
\usepackage[utf8]{inputenc}%
\usepackage{lmodern}%
\usepackage{textcomp}%
\usepackage{lastpage}%
\usepackage{geometry}%
\geometry{margin=2.2cm}%
%
\usepackage{amsmath}%
\usepackage{amssymb}%
\usepackage{booktabs}%
\usepackage{longtable}%
\usepackage{array}%
\usepackage{hyperref}%
\hypersetup{hidelinks}%
\title{IME 2026 {-} Física (claude\_haiku\_5\_5) x Poliedro\\[4pt]\large Comparação com o gabarito Poliedro Resolve\thanks{Gerado com gpt-resolve: \url{https://github.com/lgabs/gpt-resolve}}}%
\date{10/11/2026}%
%
\begin{document}%
\normalsize%
\maketitle%
\section*{Resumo}\textbf{Nota média: 10,0/10} (26 de 26 itens com a resolução correta; cada item vale o mesmo dentro da questão).\par\medskip%
\begin{longtable}{@{}c c p{0.78\textwidth}@{}}\toprule Questão & Nota & Resumo \\ \midrule
1 & 10,0 & A resolução está correta e coincide com o gabarito do Poliedro nos três itens. No item (a), os dois dão {-}2 m < x < 4 m. No item (b), os dois dão E\_sup = 32/3 J $\approx$ 10,7 J. No item (c), o esboço v$\times$t tem a mesma forma nos dois: começa em {-}v0, anula{-}se no retorno em x < {-}2, passa por +v0 em x = 0 em T2 = 2T1, tem um mínimo positivo em x = 4 e depois cresce sem limite. Refiz as contas e a integração numérica, e elas confirmam o exemplo numérico da resolução e a constante do período assintótico. Não encontrei erros na resolução. No gabarito há só um provável erro de digitação na conta do item (b), que não altera o resultado. \\
2 & 10,0 & A resolução está correta e bate com o gabarito do Poliedro nos 5 itens: a) 16/9; b) 3/4; c) M\_B/M\_CA = 1/8; d) v\_B/v\_CA = 2; e) apoastro. Conferi as contas com frações exatas em Python e a transcrição com o enunciado original. No item e, a resolução vai além do gabarito: calcula a = 9r\_A/14 e e = 5/9 e justifica por que o ponto do choque é um apsis. Não há erro em nenhum dos dois; as dúvidas são de hipótese implícita e de notação. \\
3 & 10,0 & A resolução e o gabarito Poliedro concordam nos dois itens: a) d\_min = 3?/5; b) $\rho$ = ?/14. Refiz as contas e confirmei as duas respostas: tração nula no topo do círculo de raio ? - d; folga em cos$\theta$ = 1/2 com v$^2$ = 2g?/7; e $\rho$ = (v/2)$^2$/g no apogeu do lançamento oblíquo. A resolução também justifica, com contas corretas, que o apogeu (?/28 abaixo de A) é o ponto mais alto da trajetória depois de B. Nos pontos de dúvida não há erro que afete a nota, só a idealização da barra como mero pivô e um erro de digitação no gabarito. \\
4 & 10,0 & A resolução está de acordo com o gabarito Poliedro: E = 7560 J, com malha de 2 $\Omega$ e fontes equivalentes de 40 V com 6, 3 e 2 $\Omega$, que dão 50, 128 e 200 W em cada intervalo de 20 s. A figura original confirma a topologia (R2 em série antes da primeira derivação) e a ordem de entrada dos ramos. Conferi tudo numericamente. A resolução é mais rigorosa que o gabarito: prova que a recursão converge, descarta a raiz negativa e confere o balanço de potência e a máxima transferência em k = 3. No gabarito só há falhas de digitação sem efeito no resultado (rótulo E\_1 repetido). Nota integral. \\
5 & 10,0 & A resolução e o gabarito do Poliedro chegam aos mesmos resultados nos dois itens, e conferi as contas de novo. No item (a): R\_C = 22/3 N, R\_D = 200/3 N, R\_E = 26 N e R\_G = 26 N. No item (b): m\_máx = 8,4 kg, quando R\_C = 0 e CG fica na iminência de girar em torno de D, com R\_G = 48 N. A resolução é mais completa: dá as reações em função de m e prova que não há equilíbrio para m > 8,4 kg. Não há erro de física nem de conta. A única ressalva é que o enunciado do item (b) não dá a resistência da articulação, e os dois resolveram como se o limite viesse do apoio C. \\
6 & 10,0 & A resolução está correta e coincide com o gabarito Poliedro nos três itens: a) 1,25 A; b) 56,25 J; c) 2,8125 N (45/16 N) para a direita. As duas usam o mesmo modelo: $\varepsilon$ = 10 V, r = 2 $\Omega$ e i = ($\varepsilon$ + BLv)/(R + r). As contas foram refeitas com python e conferem. A resolução também justifica o sentido da fem pela figura e fecha o balanço de energia. Só há ressalvas menores: a leitura da perspectiva da Fig. 1 decide os itens b e c, mas está correta, e há um erro de digitação sem efeito ('1 $\leq$ i'). \\
7 & 10,0 & A resolução e o gabarito do Poliedro concordam nos três itens: a) V(a) = kQ/(2a); b) V(a) = 5kQ/(6a) e V(2a) = V(3a) = kQ/(3a), com o mesmo gráfico (hipérbole kQ/r - kQ/(6a) até 2a, patamar no condutor e kQ/r fora); c) -Q nas superfícies internas, +4Q/5 na externa da casca menor e +6Q/5 na da maior. Refiz todas as contas e confirmei. A resolução está correta e mais completa, porque declara as hipóteses do fio ideal e da influência mútua desprezível e confere o item c pela energia. Não há divergência. \\
8 & 10,0 & A resolução está correta e concorda com o gabarito Poliedro nos dois itens. No item a, Q = 4,8$\times$10\^{}8 J, Q $\approx$ 83,3 kW, COP = 3/4 $\times$ 4 = 3 e P $\approx$ 27,8 kW $\approx$ 37 HP. No item b, a estimativa de 30 HP do engenheiro não procede. As contas foram refeitas em python e conferem, e o enunciado transcrito bate com o PDF original. A resolução acrescenta análises corretas na parte b: 81\% da carga, cerca de 118,5 min de processo e o caso ideal de 27,8 HP. Não há divergência real. As dúvidas listadas são apenas hipóteses implícitas e detalhes de notação. \\
9 & 10,0 & A resolução e o gabarito do Poliedro coincidem nos três itens. a) I = I\_1($\theta$){[}4 + 2$\surd$3 cos(2$\pi$a sen$\theta$/$\lambda$){]}. b) I\_máx = (4+2$\surd$3)I\_1, em sen$\theta$ = $\lambda$/a. c) I\_mín = (4-2$\surd$3)I\_1, em sen$\theta$ = $\lambda$/(2a). Conferi a transcrição com o PDF original e refiz as contas, inclusive as checagens numéricas da resolução (posição dos extremos com I\_1 = cos$^2$$\theta$ e caminho exato, razão $\approx$ 13,5, máximo espúrio para a = 0,8$\lambda$). Está tudo correto. A resolução vai além do gabarito ao declarar as condições de existência e discutir como a dependência de I\_1 em $\theta$ desloca os extremos. Nada disso altera a nota. Nota esperada: 3/3 itens corretos. \\
10 & 10,0 & A resolução e o gabarito Poliedro concordam nos dois itens. Em a), tan $\theta$ = 2$\cdot$sen($\pi$/n)/(2$\cdot$cos($\pi$/n) - $\surd$3). Em b), n $\geq$ 7 (n > 6). Refiz a conta: a lei dos senos no triângulo O{-}A{-}B com ângulo central $\pi$/n confirma a fórmula. A simulação do traçado do raio para n = 7 e n = 8 mostra o fechamento exato após 2n reflexões, com n em cada espelho. A resolução é mais completa que o gabarito: justifica que $\theta$ < $\pi$/2, mostra que existe solução para todo n $\geq$ 7 e explica por que n = 6 (feixe rasante) fica de fora. Não há erro em nenhum dos dois. Restam apenas pontos menores para levar ao professor: o caso{-}limite n = 6 e o fato de a figura ser esquemática. \\
\bottomrule\end{longtable}%
\section*{Questão 1}%
\href{https://poliedroresolve.sistemapoliedro.com.br/vestibulares/ime/2026/ime-2-fase-dia-2-24-10-2025-2-fase-dia-2/dissertativa-1-fisica-2-fase-dia-2-ime-2026}{Resolução do gabarito}\par\smallskip%
\begin{longtable}{@{}p{0.06\textwidth} p{0.33\textwidth} p{0.33\textwidth} p{0.1\textwidth} p{0.08\textwidth}@{}}\toprule Item & Resolução & Gabarito & Veredito & Certo \\ \midrule
a & {-}2 m < x < 4 m, com movimento oscilatório para 0 < E < 32/3 J; os extremos ficam de fora. & {-}2 < x < 4 (m): é o poço de potencial em torno do equilíbrio estável x = 0. Em x = 4 o equilíbrio é instável, e x = {-}2 tem a mesma Ep que x = 4 (32/3 J). & igual & ambos \\
b & E\_sup = Ep(4) = 32/3 J $\approx$ 10,7 J. É supremo e não é atingido; o período diverge como ($\surd$m/2)$\cdot$ln{[}1/(E\_sup {-} E){]}. & E\_sup = 32/3 J ($\approx$ 10,7 J), igual a Ep(4) = Ep({-}2). & igual & ambos \\
c & v(0) = {-}v0, com v0 = $\surd$(2E/m), em mínimo de tangente horizontal. v cresce até 0 em t1 (retorno em x\_t < {-}2, ponto de inflexão) e chega a +v0 em t2 = 2t1 (x = 0, máximo local). Depois cai até um mínimo v4 > 0 em t3 (x = 4) e cresce sem limite, com assíntota vertical em t* finito. A curva é antissimétrica em torno de (t1, 0) e não é periódica. & Começa em v = {-}V0. Anula{-}se em T1 (retorno em x < 0), passa por +V0 em T2 = 2T1 (x = 0), cai até um mínimo positivo em T3 (x = 4) e depois cresce sem limite, porque F > 0 para x > 4. & igual & ambos \\
\bottomrule\end{longtable}%
\begin{itemize}%
\item%
\textbf{a)} Mesmo intervalo aberto. Conferi as contas: Ep(4) = 32 {-} 64/3 = 32/3 e Ep({-}2) = 8 + 8/3 = 32/3. A fatoração x$^3$ {-} 6x$^2$ + 32 = (x{-}4)$^2$(x+2) também confere. A resolução justifica melhor a parte de 'posições possíveis': para cada x0 em ({-}2, 4) existe uma energia E < 32/3 J em que x0 fica na componente limitada que contém a origem. Ela também mostra que os extremos exigiriam E $\geq$ 32/3 J.%
\item%
\textbf{b)} Mesmo valor e mesma unidade. A resolução acrescenta a análise do caso E = E\_sup (x = 4 só é alcançado em tempo infinito, então não há periodicidade) e a divergência logarítmica do período. Conferi numericamente com m = 1: t(L$\to$2) {-} (1/4)$\cdot$ln(1/$\delta$) $\to$ 2,1089 para $\delta$ = 10$^{-}$$^2$ a 10$^{-}$$^{6}$, como ela afirma.%
\item%
\textbf{c)} Os esboços têm a mesma forma qualitativa. Refiz por RK4 o exemplo numérico da resolução (m = 1 kg, E = 16 J) e ele confere: v0 = 5,66 m/s, x\_t = {-}2,39 m, t1 = 0,640 s, t2 = 1,280 s, t3 = 2,192 s, v4 = 3,27 m/s, v(3 s) = 12,6 m/s e t* $\approx$ 3,975 s ($\approx$ 3,98 s). Os pontos do gráfico TikZ também batem. A tangente horizontal em t = 0, t2 e t3 e a inflexão em t1 estão corretas, porque dv/dt = F(x)/m e d$^2$v/dt$^2$ = F'(x)$\cdot$v/m. A resolução detalha mais que o gabarito: a fuga ocorre em tempo finito, com assíntota vertical.%
\end{itemize}%
\paragraph{Pontos de dúvida}%
\begin{itemize}%
\item%
\textit{(key)} Na resolução do gabarito, a linha do item (b) foi transcrita como 'E\_sup = 2(4)\^{}2 {-} (4(4)\^{}3)/3'. O fator 4 a mais em '4(4)\^{}3' deve ser erro de digitação ou de transcrição, porque o resultado declarado, 32/3 J, está correto: 32 {-} 64/3.%
\item%
\textit{(key)} A transcrição do gráfico do gabarito diz que a curva 'cruza o eixo t em T1 com inclinação máxima'. Isso vale só localmente, no trecho até T3: T1 é ponto de inflexão. Para t > T3 a inclinação cresce sem limite. Não afeta a resposta.%
\item%
\textit{(statement)} A massa da partícula não é dada. Por isso o esboço do item (c) só pode ser qualitativo, e a resolução usou m = 1 kg e E = 16 J como exemplo ilustrativo. Também não está explícito se o intervalo do item (a) deve ser aberto ou fechado. Ambos responderam com o intervalo aberto, o que é coerente com a energia estritamente menor que E\_sup.%
\item%
\textit{(solution)} A resolução vai muito além do pedido: assintótica do período, fuga em tempo finito, prova de unicidade e simetria. Está tudo correto, mas é longa para uma prova discursiva. O cálculo assintótico do período, que conferi numericamente, é material extra e não muda as respostas.%
\end{itemize}%
\section*{Questão 2}%
\href{https://poliedroresolve.sistemapoliedro.com.br/vestibulares/ime/2026/ime-2-fase-dia-2-24-10-2025-2-fase-dia-2/dissertativa-2-fisica-2-fase-dia-2-ime-2026}{Resolução do gabarito}\par\smallskip%
\begin{longtable}{@{}p{0.06\textwidth} p{0.33\textwidth} p{0.33\textwidth} p{0.1\textwidth} p{0.08\textwidth}@{}}\toprule Item & Resolução & Gabarito & Veredito & Certo \\ \midrule
a & r\_A/r\_B = 16/9 ($\approx$ 1,78) & R\_A/R\_B = 16/9 & igual & ambos \\
b & v\_A/v\_B = 3/4 (0,75) & v\_A/v\_B = 3/4 & igual & ambos \\
c & M\_B/M\_CA = 1/8 (0,125) & m\_B/m\_CA = 1/8 & igual & ambos \\
d & v\_B/v\_CA = 2 (v\_CA = v\_B/2) & v\_B/v\_CA = 2 & igual & ambos \\
e & Apoastro: v\_CA = v\_B/2 < v\_circ(r\_A) = 3v\_B/4, e pela vis{-}viva a = 9r\_A/14 < r\_A. Como o ponto do choque é um apsis (velocidade tangente), só pode ser o apoastro (e = 5/9, periastro a 2r\_A/7). & Apoastro: v\_CA < v\_A (0,5 v\_B < 0,75 v\_B), então pela vis{-}viva a < R\_A, logo R\_A = a + c e CA está no ponto mais afastado de S. & igual & ambos \\
\bottomrule\end{longtable}%
\begin{itemize}%
\item%
\textbf{a)} Pela 3ª lei de Kepler com T\_A = 640 e T\_B = 270 dias terrestres: (r\_A/r\_B)\^{}3 = (64/27)\^{}2, logo r\_A/r\_B = 16/9. Conferi com frações exatas. A diferença entre r e R é só de notação.%
\item%
\textbf{b)} v = sqrt(GM\_S/r) dá v\_A/v\_B = sqrt(9/16) = 3/4. A resolução chega ao mesmo valor por v = 2$\pi$r/T e confere por sqrt(r\_B/r\_A).%
\item%
\textbf{c)} Razão F\_S/F\_planeta = (M\_S/M)(R/r)\^{}2. Dividindo as duas condições: (M\_B/M\_CA)(r\_B/r\_A)\^{}2 = 81/2048, logo M\_B/M\_CA = (81/2048)(256/81) = 1/8. As duas usam CA à distância r\_A de S e o mesmo raio R. Conferido em Python.%
\item%
\textbf{d)} De E\_CA = 7E\_B com E\_B = {-}M\_B v\_B$^2$/2, M\_B = M\_CA/8 e GM\_S/r\_A = (9/16)v\_B$^2$, vem v\_CA$^2$/2 = (9/16 {-} 7/16)v\_B$^2$, ou seja, v\_CA$^2$ = v\_B$^2$/4. É o mesmo caminho do gabarito.%
\item%
\textbf{e)} Mesma conclusão e mesma justificativa central (velocidade abaixo da circular, logo a < r\_A num apsis, logo apoastro). A resolução calcula a mais: a = 9r\_A/14, e$^2$ = 25/81, e r\_p = 2r\_A/7. Verifiquei: r\_A = a(1+e) = (9/14)(14/9)r\_A. Ela também justifica que o ponto é um apsis, porque a colisão frontal deixa a velocidade de CA tangente. O gabarito só supõe isso ('se a colisão ocorre em um extremo').%
\end{itemize}%
\paragraph{Pontos de dúvida}%
\begin{itemize}%
\item%
\textit{(statement)} Item e depende de o ponto do choque ser um apsis. Isso vem de 'colisão frontal' com corpos 'em sentido contrário', o que deixa as velocidades colineares com a de A, que é tangente. O enunciado não diz isso de forma explícita. A resolução declara a hipótese; o gabarito a supõe sem justificar ('Se a colisão ocorre em um extremo da órbita...').%
\item%
\textit{(statement)} Se o momento do cometa superasse o de A, CA poderia sair em sentido retrógrado. O resultado não muda, porque energia, módulo da velocidade e condição de apsis dependem só de |v\_CA|. Nenhuma das duas resoluções comenta esse caso.%
\item%
\textit{(key)} No item c, o gabarito usa R\_A e R\_B para os raios das órbitas e d para o raio do planeta, que é o inverso da convenção usual. É só notação, mas pode confundir quem compara as duas soluções.%
\item%
\textit{(key)} A 'Solução 2' do item e do gabarito usa um argumento de referencial girante com força centrífuga. Ele é qualitativo e só mostra que CA começa a se aproximar de S, o que vale para um apsis com v < v\_circ. Está correto, mas é menos rigoroso que a Solução 1 (vis{-}viva).%
\end{itemize}%
\section*{Questão 3}%
\href{https://poliedroresolve.sistemapoliedro.com.br/vestibulares/ime/2026/ime-2-fase-dia-2-24-10-2025-2-fase-dia-2/dissertativa-3-fisica-2-fase-dia-2-ime-2026}{Resolução do gabarito}\par\smallskip%
\begin{longtable}{@{}p{0.06\textwidth} p{0.33\textwidth} p{0.33\textwidth} p{0.1\textwidth} p{0.08\textwidth}@{}}\toprule Item & Resolução & Gabarito & Veredito & Certo \\ \midrule
a & d\_min = 3?/5 (válido para 3?/5 $\leq$ d < ?), com T = 0 no topo do semicírculo de raio R = ? - d & d\_mín = 3?/5 (tração nula no topo: v$^2$ = g(? - d) e v$^2$ = 2g(2d - ?)) & igual & ambos \\
b & $\rho$ = ?/14 (o fio folga em sin$\varphi$ = 1/2, com v\_s$^2$ = 2g?/7; o apogeu do lançamento oblíquo fica ?/28 abaixo de A e lá $\rho$ = (v\_s/2)$^2$/g) & $\rho$ = ?/14 (r = 4?/7, o fio afrouxa em cos$\theta$ = 1/2, v$^2$ = 2g?/7; no topo do lançamento g$\rho$ = v$^2$cos$^2$$\theta$) & igual & ambos \\
\bottomrule\end{longtable}%
\begin{itemize}%
\item%
\textbf{a)} Refeito: v$^2$ no ponto mais baixo = 2g?; no topo do círculo de raio r = ? - d, v$^2$ = 2g? - 4gr, e a condição v$^2$ $\geq$ gr dá r $\leq$ 2?/5, logo d $\geq$ 3?/5. Conferido numericamente. A resolução chega ao mesmo valor por T($\varphi$) = (mg/R)(2d - 3R sin$\varphi$) $\geq$ 0 no topo, um caminho equivalente ao do gabarito.%
\item%
\textbf{b)} Refeito com d = 3?/7 e R = 4?/7: T = 0 quando a altura acima de B é R/2 (cos$\theta$ = 1/2 com $\theta$ medido da vertical, ou sin$\varphi$ = 1/2 na convenção da resolução; é o mesmo ponto). Daí v$^2$ = 2g?/7, com a velocidade a 60$^\circ$ da horizontal, de modo que v\_x = v/2 e $\rho$ = v\_x$^2$/g = ?/14. Conferi numericamente o que a resolução acrescenta: o fio segue frouxo durante o voo (|r|$^2$ - R$^2$ < 0 para 0 < t < t*), a esfera fica em x < 0 e só volta a esticar o fio em (0, -R) abaixo de B, e o apogeu fica em y = -?/28 relativo a A. Esse apogeu é mais alto que o ponto de folga (-?/7) e que o pêndulo depois do retensionamento (-27?/28). Portanto o ponto mais alto é mesmo o apogeu do lançamento.%
\end{itemize}%
\paragraph{Pontos de dúvida}%
\begin{itemize}%
\item%
\textit{(statement)} Item (a): para d $\geq$ ?/2, o topo da semicircunferência em torno de B, em y = ? - 2d relativo a A, fica sobre a própria barra. Com d = 3?/5 o topo cai em -?/5, entre A e B, e a esfera chegaria ao topo exatamente sobre a barra. Tanto o gabarito quanto a resolução tratam a barra só como pivô em B e ignoram a colisão; a resolução deixa isso explícito numa observação. É idealização implícita do enunciado e não muda a resposta.%
\item%
\textit{(solution)} A resolução vai além do pedido: verifica que o fio segue frouxo no voo e analisa o retensionamento (impulso radial, perda de energia, pêndulo com cos$\theta$\_max = 15/16) para garantir que o apogeu é o ponto mais alto de toda a trajetória após B. As contas estão corretas (conferidas numericamente). O gabarito supõe isso sem justificar. Não afeta a nota.%
\item%
\textit{(key)} Erro de digitação na resolução do gabarito, item (b): falta um conectivo em '2gl = 2g$\cdot$(4l)/7$\cdot$(1+cos$\theta$) + v\^{}2 v\^{}2 = (2gl)/7(3 - 4cos$\theta$)', que deveria ser '+ v$^2$ $\Rightarrow$ v$^2$ = ...'. Não muda o resultado.%
\end{itemize}%
\section*{Questão 4}%
\href{https://poliedroresolve.sistemapoliedro.com.br/vestibulares/ime/2026/ime-2-fase-dia-2-24-10-2025-2-fase-dia-2/dissertativa-4-fisica-2-fase-dia-2-ime-2026}{Resolução do gabarito}\par\smallskip%
\begin{longtable}{@{}p{0.06\textwidth} p{0.33\textwidth} p{0.33\textwidth} p{0.1\textwidth} p{0.08\textwidth}@{}}\toprule Item & Resolução & Gabarito & Veredito & Certo \\ \midrule
Q4 --- energia total dissipada na malha infinita entre t = 0 e t = 60 s & E = 7560 J (X = 2 $\Omega$; I = 5, 8, 10 A; P = 50, 128, 200 W; E = 1000 + 2560 + 4000 J) & E = 7560 J (R\_eq = 2 $\Omega$; i = 5, 8, 10 A; P = 50, 128, 200 W; E = 1000 + 2560 + 4000 J) & igual & ambos \\
\bottomrule\end{longtable}%
\begin{itemize}%
\item%
\textbf{Q4 --- energia total dissipada na malha infinita entre t = 0 e t = 60 s)} A resposta final e todos os valores intermediários coincidem. A figura do PDF confirma a topologia usada pelos dois: logo após S0 vem um R2 em série e só depois a primeira derivação R2, então X = R2 + (R2 ? X), o que dá X$^2$ - R2$\cdot$X - R2$^2$ = 0 e X = R2(1+$\surd$5)/2 = 2 $\Omega$. Também confirma a ordem das chaves: S0 liga primeiro o ramo da direita, S1 acrescenta o do meio e S2 o da esquerda, o que dá k = 1, 2, 3 ramos ativos. Refiz as contas em python. A recursão converge para 2,0000 $\Omega$, com X5 $\approx$ 2,00018. Com R\_th = 6, 3 e 2 $\Omega$, as correntes são I = 5, 8 e 10 A e as potências na malha são P = 50, 128 e 200 W. Multiplicando cada uma por 20 s e somando, E = 7560 J.%
\end{itemize}%
\paragraph{Pontos de dúvida}%
\begin{itemize}%
\item%
\textit{(key)} Erro de digitação no gabarito: as energias dos três intervalos aparecem todas com o mesmo rótulo E\_1. Deveriam ser E\_1, E\_2 e E\_3. Além disso, a passagem algébrica de R = (R2?R) + R2 está escrita sem separadores ('R\_2R R\^{}2 {-} R\_2\^{}2 = R\_2R'). Nenhum dos dois problemas muda o resultado.%
\item%
\textit{(statement)} Hipóteses que o enunciado deixa para a figura: as três baterias têm a mesma polaridade (placa longa para cima), as fontes são ideais, as chaves ficam entre ramos consecutivos e não há derivação no nó logo após S0. A figura do PDF confirma todas. A resolução declara essas hipóteses explicitamente. O gabarito as usa sem comentar.%
\item%
\textit{(solution)} A nota da resolução sobre uma topologia alternativa (com derivação em N0, X = 3 - $\surd$5 $\approx$ 0,764 $\Omega$, E $\approx$ 5460 J) está numericamente certa: dá 534,3 + 1725,5 + 3200 $\approx$ 5460 J. Como a figura descarta essa leitura, a nota é só um complemento e não afeta a resposta.%
\end{itemize}%
\section*{Questão 5}%
\href{https://poliedroresolve.sistemapoliedro.com.br/vestibulares/ime/2026/ime-2-fase-dia-2-24-10-2025-2-fase-dia-2/dissertativa-5-fisica-2-fase-dia-2-ime-2026}{Resolução do gabarito}\par\smallskip%
\begin{longtable}{@{}p{0.06\textwidth} p{0.33\textwidth} p{0.33\textwidth} p{0.1\textwidth} p{0.08\textwidth}@{}}\toprule Item & Resolução & Gabarito & Veredito & Certo \\ \midrule
a & R\_C = 22/3 N $\approx$ 7,3 N; R\_D = 200/3 N $\approx$ 66,7 N; R\_E = 26 N; R\_G = 26 N (para cima sobre GE), H\_G = 0 & F\_C = 22/3 N ($\approx$ 7,33 N); F\_D = 200/3 N ($\approx$ 66,7 N); F\_E = 26 N; F\_G = 26 N (força vertical trocada na articulação) & igual & ambos \\
b & m\_máx = 8,4 kg, limitado por R\_C = 14 - 5m/3 $\geq$ 0. No limite: R\_C = 0, R\_D = 96 N, R\_E = R\_G = 48 N & m = 8,4 kg (iminência de CG tombar em torno de D, com F\_C = 0 e F\_G = 48 N) & igual & ambos \\
\bottomrule\end{longtable}%
\begin{itemize}%
\item%
\textbf{a)} Refiz as contas. Em GE, $\Sigma$M\_E dá 2R\_G = (12+40)$\cdot$1, então R\_G = 26 N, e daí R\_E = 26 N. Em CG, $\Sigma$M\_C dá 6R\_D = 48$\cdot$4 + 26$\cdot$8 = 400, então R\_D = 200/3 N, e R\_C = 48 + 26 - 200/3 = 22/3 N. Os valores coincidem com o gabarito. O H\_G = 0 que a resolução acrescenta é coerente com o gabarito, que diz que as forças trocadas na articulação são puramente verticais.%
\item%
\textbf{b)} Os dois usam o mesmo critério: CG na iminência de girar em torno de D, com F\_C = 0. Daí 48$\cdot$2 = 2F\_G, logo F\_G = 48 N. Com F\_G = (12+10m)/2 vem m = 8,4 kg. Conferi que R\_C(m) = 14 - 5m/3 se anula em m = 42/5 = 8,4 kg. A resolução também mostra que para m > 8,4 kg não há equilíbrio nem com C apoiado nem com C solto, e que R\_G > 0 para todo m. Isso só reforça a conclusão.%
\end{itemize}%
\paragraph{Pontos de dúvida}%
\begin{itemize}%
\item%
\textit{(statement)} O item (b) pede o maior m 'de forma que a articulação consiga manter a viga em equilíbrio', mas o enunciado não dá a capacidade do dente. A resolução e o gabarito entendem que o limite vem da perda de contato em C (CG tomba em torno de D). É a leitura natural, mas há outra possível: um limite de resistência da articulação, que o enunciado não fornece.%
\item%
\textit{(key)} Erro de digitação no gabarito: diz 'forças trocadas entre os trechos CD e DE', quando deveria ser 'CG e GE'. Isso não afeta as contas.%
\item%
\textit{(solution)} Há hipóteses implícitas, mas razoáveis, aceitas por ambos: o bloco é tratado como carga concentrada a 1 m de E (x = 9 m), e o dente só empurra GE para cima. A peça clara (GE) apoia{-}se no dente da peça escura (CG), como mostra o detalhe da figura.%
\end{itemize}%
\section*{Questão 6}%
\href{https://poliedroresolve.sistemapoliedro.com.br/vestibulares/ime/2026/ime-2-fase-dia-2-24-10-2025-2-fase-dia-2/dissertativa-6-fisica-2-fase-dia-2-ime-2026}{Resolução do gabarito}\par\smallskip%
\begin{longtable}{@{}p{0.06\textwidth} p{0.33\textwidth} p{0.33\textwidth} p{0.1\textwidth} p{0.08\textwidth}@{}}\toprule Item & Resolução & Gabarito & Veredito & Certo \\ \midrule
a & |i|mín = 1,25 A (em t = 0, v = -2 m/s) & i\_mín = 5/4 A = 1,25 A (v mínima = -2 m/s) & igual & ambos \\
b & E = 56,25 J (225/4 J) & E = 56,25 J & igual & ambos \\
c & F(5 s) = 45/16 N $\approx$ 2,8125 N, para a direita (mesmo sentido de v) & F(5) = 125/16 - 5 = 2,8125 N, para a direita & igual & ambos \\
\bottomrule\end{longtable}%
\begin{itemize}%
\item%
\textbf{a)} $\varepsilon$ = 10 V e r = 2 $\Omega$ (Fig. 2). Lendo a Fig. 1, a corrente sai de A, vai ao trilho de trás e atravessa a barra de trás para a frente. Nesse sentido, a fem de movimento (v$\times$B) soma{-}se à da fonte, então i = (10 + 2,5v)/4. No intervalo i fica entre 1,25 e 3,75 A e nunca se anula, por isso o mínimo de |i| ocorre em v = -2 m/s: i = 1,25 A. Conferi com python.%
\item%
\textbf{b)} Entre 2 e 4 s, v = 2 m/s, logo i = 3,75 A e E = r$\cdot$i$^2$$\cdot$$\Delta$t = 2$\cdot$3,75$^2$$\cdot$2 = 56,25 J. O balanço de energia da resolução também fecha: 18,75 J saem pelos terminais, a força F faz 37,5 J de trabalho e a barra dissipa 56,25 J.%
\item%
\textbf{c)} Em t = 5 s: v = 1 m/s, a = -1 m/s$^2$ e i = 3,125 A. A força magnética vale iLB = 7,8125 N, contrária ao movimento. Então F = ma + iLB = -5 + 7,8125 = 2,8125 N > 0, apontando para a direita. 45/16 é o mesmo número que 125/16 - 5.%
\end{itemize}%
\paragraph{Pontos de dúvida}%
\begin{itemize}%
\item%
\textit{(statement)} Os itens b e c dependem de como se lê a perspectiva da Fig. 1: qual terminal (A ou B) vai a qual trilho e se o B vertical é normal ao plano dos trilhos. O desenho original confirma a leitura da resolução e do gabarito: A vai ao trilho de trás, a seta de i sai de A, o B vertical é perpendicular ao plano dos trilhos e as fems se somam. Com a ligação invertida, b daria 6,25 J e c daria 155/16 N para a esquerda; o item a continuaria 1,25 A. O gabarito não justifica essa leitura, mas ela parece correta.%
\item%
\textit{(solution)} Erro de digitação no parágrafo 'Ordens de grandeza': está escrito '1 $\leq$ i $\leq$ 3,75 A', mas o próprio item a mostra que o mínimo é 1,25 A. Não muda nenhuma resposta.%
\item%
\textit{(solution)} A resolução supõe, sem dizer, que a característica V\_AB = 10 - 2i vale no sentido de corrente que sai de A, ou seja, que a fonte é linear (Thévenin). É a leitura padrão, e o gabarito usa a mesma.%
\end{itemize}%
\section*{Questão 7}%
\href{https://poliedroresolve.sistemapoliedro.com.br/vestibulares/ime/2026/ime-2-fase-dia-2-24-10-2025-2-fase-dia-2/dissertativa-7-fisica-2-fase-dia-2-ime-2026}{Resolução do gabarito}\par\smallskip%
\begin{longtable}{@{}p{0.06\textwidth} p{0.33\textwidth} p{0.33\textwidth} p{0.1\textwidth} p{0.08\textwidth}@{}}\toprule Item & Resolução & Gabarito & Veredito & Certo \\ \midrule
a & V(a) = kQ/(2a) & V$_{1}$(a) = kQ/(2a) & igual & ambos \\
b & V(r) = kQ/r - kQ/(6a) para 0 < r $\leq$ 2a; V = kQ/(3a) para 2a $\leq$ r $\leq$ 3a; V = kQ/r para r $\geq$ 3a. V(a) = 5kQ/(6a), V(2a) = V(3a) = kQ/(3a). O gráfico é contínuo, não crescente, com patamar no condutor e V $\to$ +$\infty$ quando r $\to$ 0, desenhado também para r < 0 (simétrico) & V$_{2}$ = kQ(1/r - 1/(6a)) para r $\leq$ 2a; V$_{2}$ = kQ/(3a) para 2a $\leq$ r $\leq$ 3a; V$_{2}$ = kQ/r para r $\geq$ 3a. V$_{2}$(a) = 5kQ/(6a), V$_{2}$(2a) = V$_{2}$(3a) = kQ/(3a) & igual & ambos \\
c & Casca 1: interna -Q, externa +4Q/5. Casca 2: interna -Q, externa +6Q/5. Potencial final comum 2kQ/(5a) & Casca da Fig. 1: interna -Q, externa +4Q/5. Casca da Fig. 2: interna -Q, externa +6Q/5 & igual & ambos \\
\bottomrule\end{longtable}%
\begin{itemize}%
\item%
\textbf{a)} Por indução, a superfície interna (r = a) fica com -Q e a externa (r = 2a) com +Q. Superpondo as três cascas/cargas: V(a) = kQ/a - kQ/a + kQ/(2a) = kQ/(2a). A resolução chega ao mesmo valor porque o metal é equipotencial (V(a) = V(2a) = kQ/(2a)) e confere por integração na cavidade. Conta verificada em python.%
\item%
\textbf{b)} As expressões por trecho e os três valores pedidos coincidem. Verificação: V(a) = kQ/a - kQ/(2a) + kQ/(3a) = 5kQ/(6a). A forma do gráfico é a mesma (hipérbole deslocada até 2a, patamar até 3a, depois kQ/r). O esboço da resolução, em unidades de a e de kQ/a, está consistente com 1/x - 1/6. A extensão para r < 0 é só uma simetrização coerente com as marcas -2a e -3a da figura.%
\item%
\textbf{c)} Gauss no metal dá -Q em cada superfície interna. A conservação dá q$_{1}$ + q$_{2}$ = 2Q, e a equipotencialidade com cascas muito afastadas dá kq$_{1}$/(2a) = kq$_{2}$/(3a). Daí q$_{1}$ = 4Q/5 e q$_{2}$ = 6Q/5, com potencial comum 2kQ/(5a). A resolução ainda confere pela minimização de energia: U cai de 5kQ$^2$/(12a) para 2kQ$^2$/(5a), uma variação de kQ$^2$/(60a). Tudo verificado em python.%
\end{itemize}%
\paragraph{Pontos de dúvida}%
\begin{itemize}%
\item%
\textit{(statement)} O enunciado não diz que o fio não retém carga ao ser retirado nem que a influência mútua entre as cascas é desprezível. Ambos (gabarito e resolução) supõem isso; só a resolução declara explicitamente. É a leitura pretendida ('cascas muito distantes'), mas a hipótese do fio sem capacitância está implícita no gabarito.%
\item%
\textit{(solution)} A resolução desenha o gráfico também para r < 0, por simetria. O gabarito mostra só r $\geq$ 0. As duas formas são aceitáveis, porque o enunciado marca -2a e -3a na figura, mas pede valores só em r = a, 2a e 3a.%
\end{itemize}%
\section*{Questão 8}%
\href{https://poliedroresolve.sistemapoliedro.com.br/vestibulares/ime/2026/ime-2-fase-dia-2-24-10-2025-2-fase-dia-2/dissertativa-8-fisica-2-fase-dia-2-ime-2026}{Resolução do gabarito}\par\smallskip%
\begin{longtable}{@{}p{0.06\textwidth} p{0.33\textwidth} p{0.33\textwidth} p{0.1\textwidth} p{0.08\textwidth}@{}}\toprule Item & Resolução & Gabarito & Veredito & Certo \\ \midrule
a & Q = 120 000 kcal = 4,8$\times$10\^{}8 J; Q\_F = 250/3 kW $\approx$ 83,3 kW; COP\_Carnot = 246/61,5 = 4; COP = 3; $\dot W$ = 250 000/9 W $\approx$ 27,8 kW $\approx$ 37 HP (exato 1000/27 HP) & Q\_f = 120 000 kcal = 480 000 kJ; Q\_f = 250/3 kW = 1000/9 HP; $\beta$ = 3; P $\approx$ 37 HP & igual & ambos \\
b & Não é procedente: 30 HP (22,5 kW) extraem só 67,5 kW, 81\% dos 83,3 kW necessários; o congelamento levaria cerca de 118,5 min em vez de 96 min; faltam cerca de 7 HP, sem folga; 37 HP é um piso & Não procede: são necessários cerca de 37 HP, mais que os 30 HP indicados & igual & ambos \\
\bottomrule\end{longtable}%
\begin{itemize}%
\item%
\textbf{a)} Os passos e os valores intermediários coincidem. Refiz as contas em python: Q = (35 000 + 72 500 + 12 500) kcal $\times$ 4000 J/kcal = 4,8$\times$10\^{}8 J; Q = 4,8$\times$10\^{}8/5760 $\approx$ 83 333 W; $\beta$\_C = 246/61,5 = 4; $\beta$ = 3; P = 83 333/3 $\approx$ 27 778 W, e 27 778/750 $\approx$ 37,04 HP. O enunciado transcrito confere com o PDF original.%
\item%
\textbf{b)} A conclusão é a mesma: a estimativa não procede, porque 37 HP > 30 HP. A resolução acrescenta números de apoio, todos conferidos: 3 $\times$ 22,5 = 67,5 kW, que é 0,81 de 83,3 kW; t = 4,8$\times$10\^{}8/67 500 $\approx$ 7111 s $\approx$ 118,5 min; o COP necessário com 30 HP seria 3,70, ou 92,6\% do Carnot; a hipótese de bomba de calor dá 29,6 HP; o caso ideal de Carnot dá 27,8 HP. Esses acréscimos estão corretos e não contradizem o gabarito.%
\end{itemize}%
\paragraph{Pontos de dúvida}%
\begin{itemize}%
\item%
\textit{(statement)} O enunciado não informa como a carga térmica varia no tempo. Tanto a resolução quanto o gabarito supõem extração de calor a taxa constante (potência média = Q/$\Delta$t) e que a câmara não recebe calor do ambiente por paredes ou portas. Essa é a leitura esperada, mas é uma hipótese implícita.%
\item%
\textit{(solution)} Na parte b, a resolução especula sobre a origem dos 30 HP (Carnot sem o fator 3/4, ou COP de bomba de calor). A própria resolução avisa que os dados não confirmam isso. É um texto extra, não exigido, e não pesa na nota.%
\item%
\textit{(key)} O gabarito escreve a variação de temperatura da fase sólida como (27 - 2), resultado 25 $^\circ$C; o correto, como na resolução, é -2 - (-27). Ele também escreve 'Q\_f\^{}? = 1000/9 HP', que é a taxa de calor expressa em HP (111,1 HP). As duas coisas são só notação e não alteram o resultado.%
\item%
\textit{(solution)} A resolução converte as temperaturas para kelvin somando 273 (246 K e 307,5 K), igual ao gabarito. Com 273,15 o COP de Carnot passaria a cerca de 4,002, diferença desprezível que a própria resolução já aponta.%
\end{itemize}%
\section*{Questão 9}%
\href{https://poliedroresolve.sistemapoliedro.com.br/vestibulares/ime/2026/ime-2-fase-dia-2-24-10-2025-2-fase-dia-2/dissertativa-9-fisica-2-fase-dia-2-ime-2026}{Resolução do gabarito}\par\smallskip%
\begin{longtable}{@{}p{0.06\textwidth} p{0.33\textwidth} p{0.33\textwidth} p{0.1\textwidth} p{0.08\textwidth}@{}}\toprule Item & Resolução & Gabarito & Veredito & Certo \\ \midrule
a & I($\theta$) = I\_1($\theta$){[}4 + 2$\surd$3$\cdot$cos(2$\pi$a$\cdot$sen$\theta$/$\lambda$){]} & I\_R($\theta$) = I\_1($\theta$){[}4 + 2$\surd$3$\cdot$cos((2$\pi$/$\lambda$)a$\cdot$sen$\theta$){]} = 2I\_1($\theta$){[}2 + $\surd$3$\cdot$cos((2$\pi$a/$\lambda$)sen$\theta$){]} & igual & ambos \\
b & sen $\theta$\_1 = $\lambda$/a; I\_máx = (4 + 2$\surd$3)$\cdot$I\_1($\theta$\_1) = (1+$\surd$3)$^2$$\cdot$I\_1($\theta$\_1) $\approx$ 7,46$\cdot$I\_1($\theta$\_1), para a > $\lambda$ & I\_1$^\circ$máx = I\_1($\theta$)(4 + 2$\surd$3), em sen$\theta$ = $\lambda$/a & igual & ambos \\
c & sen $\theta$\_2 = $\lambda$/(2a); I\_mín = (4 - 2$\surd$3)$\cdot$I\_1($\theta$\_2) = ($\surd$3-1)$^2$$\cdot$I\_1($\theta$\_2) $\approx$ 0,536$\cdot$I\_1($\theta$\_2), para a > $\lambda$/2 & I\_1$^\circ$mín = I\_1($\theta$)(4 - 2$\surd$3), em sen$\theta$ = $\lambda$/(2a) & igual & ambos \\
\bottomrule\end{longtable}%
\begin{itemize}%
\item%
\textbf{a)} A expressão é a mesma. As duas partem de I\_2 = 3I\_1, I = I\_1 + I\_2 + 2$\surd$(I\_1I\_2)cos$\Delta$$\varphi$ e $\Delta$$\varphi$ = (2$\pi$/$\lambda$)a$\cdot$sen$\theta$, com raios quase paralelos porque L ? a. O sentido de $\theta$ não altera o resultado, porque o cosseno é par. A transcrição foi conferida com a pág. 9 do PDF e está fiel.%
\item%
\textbf{b)} O fator (4+2$\surd$3) e a posição (sen$\theta$ = $\lambda$/a, ou seja, $\Delta$$\varphi$ = 2$\pi$, sem contar o máximo central em $\theta$ = 0) coincidem. Escrever I\_1($\theta$\_1) em vez de I\_1($\theta$) é só notação. A resolução acrescenta a condição de existência a > $\lambda$, que é correta e não está no gabarito. Conferi numericamente com I\_1 = cos$^2$$\theta$, a = 5$\lambda$ e diferença de caminho exata (L = 10$^{6}$$\lambda$). O máximo exato fica em 0,20043 rad, contra arcsen(0,2) = 0,20136 rad, e o valor difere de (4+2$\surd$3)I\_1($\theta$\_1) em 0,019\%. Isso confirma as conferências da própria resolução.%
\item%
\textbf{c)} O fator (4-2$\surd$3) e a posição ($\Delta$$\varphi$ = $\pi$, ou seja, sen$\theta$ = $\lambda$/(2a)) coincidem. O mínimo não se anula porque as amplitudes são diferentes (A\_2 = $\surd$3$\cdot$A\_1), e os dois tratam isso corretamente. A resolução acrescenta a condição a > $\lambda$/2. Na conferência numérica (I\_1 = cos$^2$$\theta$, a = 5$\lambda$), o mínimo exato fica em 0,10020 rad, contra 0,10017 rad previsto, com I praticamente igual a (4-2$\surd$3)I\_1($\theta$\_2).%
\end{itemize}%
\paragraph{Pontos de dúvida}%
\begin{itemize}%
\item%
\textit{(statement)} Se I\_1 depende de $\theta$, os máximos e mínimos locais de I($\theta$) = I\_1($\theta$)$\cdot$g($\theta$) não ficam exatamente em $\Delta$$\varphi$ = 2$\pi$ e $\Delta$$\varphi$ = $\pi$. Eles se deslocam um pouco, da ordem de |d ln I\_1/d$\theta$|$\cdot$$\lambda$$^2$/a$^2$. Para a próximo de $\lambda$, a variação de I\_1 pode até criar extremos que não vêm da interferência: com I\_1 = cos$^2$$\theta$ e a = 0,8$\lambda$, aparece um máximo local perto de 1,02 rad, o que conferi numericamente. O gabarito ignora isso e a resolução comenta. Como o enunciado pede a resposta 'em função de I\_1($\theta$)', a leitura pretendida é a dos extremos do fator de interferência, e os dois seguem essa leitura.%
\item%
\textit{(key)} Nos itens b e c, o gabarito escreve I\_1($\theta$) sem dizer que se trata do valor no ângulo do extremo ($\theta$ = arcsen($\lambda$/a) ou arcsen($\lambda$/(2a))). A resolução escreve I\_1($\theta$\_1) e I\_1($\theta$\_2). A diferença é só de notação.%
\item%
\textit{(key)} O gabarito não menciona as condições de existência: a > $\lambda$ para o primeiro máximo com $\theta$ > 0 e a > $\lambda$/2 para o primeiro mínimo. A resolução as declara. Não muda a resposta, mas completa a solução.%
\end{itemize}%
\section*{Questão 10}%
\href{https://poliedroresolve.sistemapoliedro.com.br/vestibulares/ime/2026/ime-2-fase-dia-2-24-10-2025-2-fase-dia-2/dissertativa-10-fisica-2-fase-dia-2-ime-2026}{Resolução do gabarito}\par\smallskip%
\begin{longtable}{@{}p{0.06\textwidth} p{0.33\textwidth} p{0.33\textwidth} p{0.1\textwidth} p{0.08\textwidth}@{}}\toprule Item & Resolução & Gabarito & Veredito & Certo \\ \midrule
a & tan $\theta$ = sen($\pi$/n)/(cos($\pi$/n) - $\surd$3/2) = 2$\cdot$sen($\pi$/n)/(2$\cdot$cos($\pi$/n) - $\surd$3), para n $\neq$ 6 & tg $\theta$ = 2$\cdot$sen($\pi$/n)/(2$\cdot$cos($\pi$/n) - $\surd$3) & igual & ambos \\
b & n $\geq$ 7, ou seja, n $\in$ \{7, 8, 9, \ldots{}\}. Para n $\leq$ 5, tan $\theta$ < 0 (e n $\leq$ 2 já fica fora porque $\pi$/n < $\theta$ < $\pi$/2). n = 6 é rasante ($\theta$ = 90$^\circ$) e foi excluído. & n > 6 (n natural, ou seja, n $\geq$ 7), porque tg $\theta$ > 0 exige cos($\pi$/n) > $\surd$3/2 & igual & ambos \\
\bottomrule\end{longtable}%
\begin{itemize}%
\item%
\textbf{a)} A expressão é a mesma. As duas partem do triângulo O{-}A{-}B. O ângulo central entre reflexões consecutivas é $\varphi$ = $\pi$/n, porque são 2n reflexões por volta. Pela lei dos senos, R/sen $\theta$ = r/sen $\alpha$, com $\alpha$ = $\theta$ - $\varphi$. Isso dá sen($\theta$ - $\pi$/n) = ($\surd$3/2)$\cdot$sen $\theta$ e, daí, a fórmula. Conferi em Python que, com esse $\theta$, a lei dos senos fica satisfeita até o erro de máquina (n = 7, 8, 10, 20). Também simulei o traçado do raio para n = 7 e n = 8: o feixe alterna entre os raios R e r, os pontos ficam igualmente espaçados de $\pi$/n e, após 2n reflexões, volta exatamente ao ponto A (n pontos em cada espelho). A ressalva n $\neq$ 6 da resolução é coerente, pois aí o denominador se anula.%
\item%
\textbf{b)} Mesma conclusão. Como 0 < $\theta$ < $\pi$/2 e sen($\pi$/n) > 0, é preciso cos($\pi$/n) > cos($\pi$/6), o que dá n > 6. Conferi numericamente: n = 3, 4 e 5 dão tan $\theta$ negativo (-2,37; -4,45; -10,31). Para n = 7, tan $\theta$ $\approx$ 12,42 ($\theta$ $\approx$ 85,40$^\circ$). Para n = 8, tan $\theta$ $\approx$ 6,61 ($\theta$ $\approx$ 81,40$^\circ$). $\theta$ diminui à medida que n cresce. A resolução vai além do gabarito: justifica por que $\theta$ < $\pi$/2 (o feixe precisa se afastar do espelho interno), mostra que existe solução para todo n $\geq$ 7 e explica a exclusão de n = 6 (o feixe só tangencia o cilindro interno e a trajetória vira um hexágono inscrito). O gabarito exclui n = 6 apenas pela desigualdade estrita.%
\end{itemize}%
\paragraph{Pontos de dúvida}%
\begin{itemize}%
\item%
\textit{(key)} Em b, o gabarito usa 0 < $\theta$ < $\pi$/2 sem justificar e exclui n = 6 só pela desigualdade estrita. Em n = 6 o feixe tangencia o cilindro interno ($\theta$ = 90$^\circ$), sem desvio. Um aluno poderia discutir se esse caso{-}limite conta como 'ponto de reflexão'. Tanto a resolução quanto o gabarito o excluem, o que é a leitura mais razoável.%
\item%
\textit{(statement)} A figura do enunciado é esquemática. Nela, r/R $\approx$ 0,57 e $\theta$ $\approx$ 45$^\circ$, o que não combina com r/R = $\surd$3/2. Pela fórmula, $\theta$ $\approx$ 45$^\circ$ só ocorreria com n $\approx$ 25, e para n = 7 o ângulo seria $\theta$ $\approx$ 85$^\circ$. Isso não afeta a resposta, mas pode confundir quem tente tirar valores do desenho.%
\item%
\textit{(statement)} A interpretação de 'os raios se sobrepõem após uma volta' como 2n$\cdot$$\Delta$$\varphi$ = 2$\pi$, com n reflexões em cada espelho e retorno ao ponto A na primeira volta, é uma hipótese comum à resolução e ao gabarito. Uma leitura com mais de uma volta mudaria a condição de fechamento, mas o enunciado fala em 'uma volta'.%
\end{itemize}%
\end{document}